\documentclass[11pt,a4paper]{article}
\usepackage[top=0.75in,bottom=0.75in,left=0.75in,right=0.75in]{geometry}
\usepackage[T1]{fontenc}
\usepackage{times,microtype,enumitem,hyperref,titlesec}
\pagestyle{empty}
\setlength{\parindent}{0pt}
\setlength{\parskip}{0pt}
\linespread{1.10}
\hypersetup{hidelinks,breaklinks=true}
\titleformat{\section}{\bfseries\normalsize}{}{0pt}{}[\titlerule]
\titlespacing*{\section}{0pt}{10pt}{4pt}
\newcommand{\entry}[3]{\noindent\textbf{#1}\hfill\textbf{#3}\\[-1pt]\textit{#2}\par}
\newenvironment{cvitems}{\begin{itemize}[leftmargin=1.1em,label=\textbullet,itemsep=2pt,topsep=3pt,parsep=0pt,partopsep=0pt]}{\end{itemize}}
\begin{document}
{\LARGE\bfseries Aarush Ram Anandh}\par\vspace{4pt}
\href{mailto:aarush.anandh_ug2024@ashoka.edu.in}{aarush.anandh\_ug2024@ashoka.edu.in}\quad|\quad +91 9659786783\par\vspace{2pt}
\href{https://github.com/MushuRam01}{github.com/MushuRam01}\quad|\quad
\href{https://www.linkedin.com/in/aarush-ram-anandh/}{linkedin.com/in/aarush-ram-anandh}\par
\vspace{20pt}
12 November 2026\par\vspace{12pt}
Hiring Team\\
Google\par\vspace{12pt}
\textbf{Application: Student Researcher}\par\vspace{14pt}
Dear Hiring Team,\par\vspace{10pt}
I am applying to Google's Student Researcher program for consideration across research teams. I am completing a B.Sc. (Hons.) in Computer Science and Mathematics at Ashoka University, with graduation expected in 2028. My work connects mathematical questions about models with the practical work of implementing methods, designing experiments and understanding their results. I would bring that combination to an exploratory research placement, with interests spanning learning algorithms, generative models and data-driven scientific problems.\par\vspace{10pt}
At the Constrained Image-Synthesis Lab, I investigated how to preserve specified colours during latent-diffusion sampling. The work achieved $\Delta E_{00}<1$ and approximately four times better colour accuracy than evaluated open-source baselines, and is reported in our public preprint, \textit{Inference-Time Loss-Guided Colour Preservation in Diffusion Sampling} (arXiv:2601.17259). My ongoing Controlled Generalised Sliced Wasserstein research extends this interest in controllable methods to optimal transport, connecting mathematical properties, C++ implementations, Lean specifications/proofs and a hyperspectral experiment. A paper on this work has been submitted to JMLR.\par\vspace{10pt}
I also bring experience examining what an experiment establishes. In a two-person RL project, we evaluated PPO across two team variants and five seeds each. The agents achieved 99.7--100\% wins against a fixed heuristic; statistical tests found no significant learning-speed difference between variants. At the Koita Centre for Digital Health, I contributed to a wearable-data pipeline achieving approximately 70\% sensitivity for asthma exacerbation detection in study evaluation. Work on 191 patients involved signal processing, patient-wise splits and bootstrap intervals, connecting modelling decisions to evaluation on noisy longitudinal data.\par\vspace{10pt}
My industry experience adds practice improving systems under concrete constraints. At Avartan Labs, I raised rare-class correct detection from 40\% to 80\% in an RT-DETR sorting pipeline covering 18 material classes, alongside failure analysis and targeted model improvements. This complements my academic work with experience turning model diagnostics into a measurable improvement on an applied task.\par\vspace{10pt}
Across these settings, I can contribute Python/PyTorch and C++ implementation, mathematical reasoning, collaborative experimentation and clear communication of findings and limitations. Google's program-wide matching process appeals to me because these skills can support different research questions. I would welcome the opportunity to develop them further while contributing to a team's scientific and engineering work. Thank you for considering my application.\par\vspace{14pt}
Yours sincerely,\\[5pt]
Aarush Ram Anandh
\end{document}
